\lyseisE{Επανάληψη κεφαλαίου Γ2}
\solr{Γ2}{1} $\dfrac{7+9+12+4+8}{5}=\dfrac{40}{5}=8$.
\solr{Γ2}{2} Στη σειρά: $3,6,8,11,15$, διάμεσος $8$. Στη σειρά: $2,4,7,9$, διάμεσος $\dfrac{4+7}{2}=5{,}5$.
\solr{Γ2}{3} $12\cdot6{,}5=78$.
\solr{Γ2}{4} Θέλουμε άθροισμα $5\cdot15=75$. Έως τώρα $13+16+14+15=58$. Χρειάζεται $75-58=17$.
\solr{Γ2}{5} Γινόμενα: $5,16,12,12$, άθροισμα $45$. Σύνολο $5+8+4+3=20$. Μέση τιμή $45:20=2{,}25$ ώρες.
\solr{Γ2}{6} Κέντρα: $10,30,50$. Γινόμενα: $40,300,300$, άθροισμα $640$. Μέση τιμή $640:20=32$ λεπτά.
\solr{Γ2}{7} Εύρος $22-14=8$ έτη. Το $14$ εμφανίζεται τρεις φορές, άρα κορυφή το $14$.
\solr{Γ2}{8} Απόμακρη τιμή το $2$. Με αυτή: $\dfrac{78}{6}=13$. Χωρίς αυτή: $\dfrac{76}{5}=15{,}2$.
\solr{Γ2}{9} Πρώτος: μέση τιμή $\dfrac{40}{5}=8$, εύρος $14-2=12$. Δεύτερος: μέση τιμή $\dfrac{40}{5}=8$, εύρος $9-7=2$. Ίδιος μέσος όρος, αλλά ο δεύτερος είναι πολύ πιο σταθερός, ο πρώτος πολύ πιο απρόβλεπτος.
\solr{Γ2}{10} Εύρος $10-0=10$ ώρες. Κορυφή η κλάση $4$--$6$ (ψηλότερο ορθογώνιο, $7$ μαθητές). Κέντρα $1,3,5,7,9$: άθροισμα $1\cdot1+3\cdot4+5\cdot7+7\cdot5+9\cdot3=1+12+35+35+27=110$, μέση τιμή $110:20=5{,}5$ ώρες.
\solr{Γ2}{11} Το αρχικό άθροισμα είναι $10\cdot8=80$. Νέο άθροισμα $80+19=99$, νέο πλήθος $11$. Νέα μέση τιμή $99:11=9$.
\solr{Γ2}{12} Άθροισμα βαθμών $5\cdot12+15\cdot16=60+240=300$. Σύνολο μαθητών $20$. Μέσος όρος $300:20=15$.
\solr{Γ2}{13} Παράδειγμα: $2,4,5,6,9,10$. Διάμεσος $\dfrac{5+6}{2}=5{,}5$, εύρος $10-2=8$, άθροισμα $36$ άρα μέση τιμή $6$.
\solr{Γ2}{14} Άθροισμα $4$, μέση τιμή $4:6\approx0{,}67$ $^\circ$C. Στη σειρά: $-3,-1,0,2,2,4$, διάμεσος $\dfrac{0+2}{2}=1$. Εύρος $4-(-3)=7$. Κορυφή το $2$.
\solr{Γ2}{15} α) Λάθος (η μεσαία τιμή). β) Σωστό. γ) Σωστό ($3$ και $5$). δ) Λάθος. ε) Λάθος: $4-(-5)=9$.
\solr{Γ2}{16} α) Λάθος. β) Σωστό. γ) Σωστό. δ) Λάθος: το εύρος μένει ίδιο. ε) Σωστό.
